\clearpage\begingroup\SourceLeaf{44b}\HistoricalHead{39}{R}{GERLING AN GAUSS.}
\begin{center}
[6.]\\[.7em]
\textsc{Gauss} an \textsc{Gerling}. Göttingen, 30. December 1852.
\end{center}
\noindent .\ .\ .\ .\ .\ .\ .\ .\ .\ .\ .\ .\ .\ .\ .\ .\ .\ .\ .\ .\ .
\par\noindent\hspace*{\parindent}Den Apparat für den \textsc{Foucauld}schen Versuch, wovon ich Ihnen früher
geschrieben habe, denke ich ausführen zu lassen. Ich hoffe, dass dadurch
das Phänomen in jedem Local allemahl schon nach sehr kurzer Zeit bestimmt
hervortretend gemacht werden kann. Die Fallversuche nach \textsc{Guglielmini}[\textsuperscript{*})]\begingroup\renewcommand{\thefootnote}{}\footnotetext{[\textsuperscript{*}) G. B. \textsc{Guglielmini}, \textit{De diurno terrae motu, experimentis physico-mathematicis confirmato opusculum}, Bononiae 1792.]}\endgroup\
u. a. sind eigentlich wenig geeignet, die Drehungsbewegungen der Erde erkennbar
zu machen, da sie nach den kostspieligsten Zurüstungen doch immer
nur höchst rohe Resultate geben können. Die Versuche von \textsc{Hook}[\textsuperscript{**})]\begingroup\renewcommand{\thefootnote}{}\footnotetext{[\textsuperscript{**}) \textsc{Robert Hooke}, \textit{An attempt to prove the Motion of the Earth}, London 1674.]}\endgroup\ und
selbst die von \textsc{Guglielmini} beweisen eigentlich gar nichts, da bei letzteren
das Loth erst \(\tfrac{1}{2}\) Jahr nach den Fallversuchen angewandt wurde. Wenn Sie
übrigens sagen, alle Versuche hätten eine Bewegung nach Süden ergeben,
wofür die Theorie keine Rechenschaft hätte, so verstehe ich diess nicht.
\textsc{Benzenbergs} Versuche in Hamburg gaben zwar eine solche, aber \textit{so}, dass
dieses Resultat gar keinen Werth hatte, wie er selbst \textit{nicht} verkannte, und
deshalb unternahm er ja eben die neuen Versuche in einem Schacht im
Bergischen und diese gaben keine Abweichung nach Süden, sondern nur
nach \textls[120]{Norden} (s. sein Buch[\textsuperscript{***})]\begingroup\renewcommand{\thefootnote}{}\footnotetext{[\textsuperscript{***}) \textsc{Joh. Friedrich Benzenberg}, \textit{Versuche über das Gesetz des Falles, über den Widerstand der Luft und über die Umdrehung der Erde}, Dortmund 1804; die von \textsc{Gauss} erwähnten Versuche machte \textsc{Benzenberg} 1802 in Hamburg auf dem Michaelisturm und 1804 in einem Kohlenschacht bei Schlebusch in der Grafschaft Mark. Vergl. auch J. F. \textsc{Benzenberg}, \textit{Versuche über die Umdrehung der Erde, neu berechnet}, Düsseldorf 1845.]}\endgroup\ p. 425). Bei Veranlassung Ihrer obigen
\clearpage\SourceLeaf{45}\HistoricalHead{40}{L}{PHYSIK. BRIEFWECHSEL.}
\noindent Aeusserung habe ich übrigens die \textsc{Reich}schen Versuche[\textsuperscript{*})]\begingroup\renewcommand{\thefootnote}{}\footnotetext{[\textsuperscript{*}) \textsc{Ferdinand Reich}, \textit{Fallversuche über die Umdrehung der Erde u.s.w.}, Freiberg 1832.]}\endgroup\ wieder angesehen.
Derselbe findet nach \textls[120]{seiner} Rechnung (p. 46) eine Abweichung nach Süden
von \(4{,}374\) Millimeter mit einem wahrscheinlichen Fehler von \(2^{\mathrm{mm}}{,}700\). Auch
diess als richtig angenommen, würde ich doch seine folgenden Zeilen »Was
die letztere betrifft .\ .\ .\ .\ .\ .\ .\ nach Süden« nicht billigen können. Von einem
Resultate, dessen absolute Grösse nur \(1\tfrac{1}{2}\)mahl so gross ist als der wahrscheinliche
Fehler, darf man nicht sagen, dass sie »noch nicht ausser allen Zweifel
gesetzt sei«, sondern nur 1) in dem Fall, wo man gar nichts weiter davon
weiss, als den Ausfall der Versuche, — dass daraus für die Realität der Grösse
nur erst eine mässige Wahrscheinlichkeit resultire, hingegen 2) in dem Fall,
wo \textls[120]{andere} gewichtige Gründe gegen die Realität sprechen, würde ich sagen,
dass auf den Ausfall nur wenig zu geben sei, und 3) wenn wie hier seit 50
Jahren durch strenge Theorie \textls[120]{bewiesen} ist, dass eine Abweichung nach Süden
gar nicht stattnehmig sei, würde man sagen müssen, dass gar \textls[120]{nichts} darauf
zu geben sei. Was den letzten Theil von \textsc{Reichs} Phrasen betrifft, so finde
ich die Sache ungefähr eben so sonderbar, wie den Umstand, dass alle 3 jetzigen
Kaiser am 2. December ihren Thron bestiegen haben, und zwar \textsc{Franz Josef}
4 Jahr, \textsc{Nicolaus} 27 Jahre früher als \textsc{Nap[oleon]} III (ein Mathematiker könnte
die Sonderbarkeit noch erhöhen, wenn er die Zahlen so schriebe \(4=2^2\),
\(27=3^3\)). Mit solchen Phantasiespielen belustigt man sich wohl, aber niemand
legt ihnen eine ernsthafte Bedeutung bei.

Nun aber kommt noch ein Hauptpunkt. Ich finde \textsc{Reichs} Rechnung
\textls[120]{falsch}. Wie er es angefangen hat, obige Zahlen herauszubringen, weiss ich
nicht: aber eine richtige Rechnung giebt anstatt der Zahlen p. 45
\begin{center}
\begin{tabular}{rr}
28,282 & 2,703\\
5,061 & 2,700
\end{tabular}
\end{center}
\noindent folgende:
\par\smallskip\noindent\hspace*{\parindent}Hauptresultat für die \textls[120]{östliche} Abweichung
\begin{center}
\(28^{\mathrm{mm}}{,}527\),\quad wahrscheinlicher Fehler\quad \(3{,}481\),
\end{center}
\noindent\hspace*{\parindent}Hauptresultat für die südliche Abweichung
\begin{center}
\(3^{\mathrm{mm}}{,}540\),\quad [wahrscheinlicher Fehler]\quad \(3{,}400\).
\end{center}
\clearpage\SourceLeaf{46}\HistoricalHead{41}{R}{GAUSS AN GERLING.}
\noindent Nach Correction der falschen Orientirung erhalte ich hieraus:
\begin{center}
\begin{tabular}{@{}lr@{\,}r@{}l@{}l@{}}
Abweichung nach & Osten & 28 & \({}^{\mathrm{mm}}\) & ,605,\\
 & Süden & 2 & \(\phantom{{}^{\mathrm{mm}}}\) & ,874,
\end{tabular}
\end{center}
\noindent also die vermeintliche Abweichung nach Süden noch kleiner als der wahrscheinliche
Fehler darin.

Wenn Ihrem Briefe zufolge Ihr »Desiderium wegen des \textsc{Foucauld}schen
Versuches noch immer fortbesteht«, so kann ich darauf nur erwiedern, dass
meiner Ueberzeugung nach der Grund der für Sie Statt findenden Dunkelheit
darin liegt, dass Sie die Aufgabe für leichter ansehen, als sie ist, und erwarten,
dass sie sich durch ein Aperçu (um ein Wort zu gebrauchen, was
\textsc{Göthe} so oft zu seinem cheval de bataille machte) beantworten lasse, zu
welcher Erwartung Sie gar kein Recht haben. Es gibt sehr viele Fälle, wo
sich durch ein Aperçu in eine intricate Frage ein helles Licht bringen lässt,
aber auch andere, wo dies Licht nur ein Irrlicht ist, und höchstens die Wirkung
hat, dass mancher sich einbildet, nun die Sache zu verstehen, obwohl
er eigentlich sich darin bloss täuscht. Ich rechne dahin den Versuch des
p. \textsc{Schaub} in den A. N. Nro. 838[\textsuperscript{*})]\begingroup\renewcommand{\thefootnote}{}\footnotetext{[\textsuperscript{*}) F. \textsc{Schaub}, \textit{Elementarer Beweis der Wirkung der Umdrehung der Erde auf die Schwingungsebene des Pendels}, Astronomische Nachrichten \textbf{35}, 1853, Spalte 353.]}\endgroup, der kein elementarer Beweis ist, sondern
gar kein Beweis. Der Hergang der Sache im absoluten Raume kann so gedacht
werden. Ein materieller Punkt \(A\) wird nach einem festen Punkte \(C\) zu
angezogen, zugleich aber ist er genöthigt, von einem dritten Punkt \(B\) immer
in einer unveränderlichen Entfernung zu bleiben: man sucht die Bewegung
des Punktes \(A\) im Raum. Ist nun \(B\) gleichfalls fest und nimmt [man] die
Stärke der Anziehung nach \(C\) zu wie constant an, auch überhaupt \(AB\) wie
verschwindend gegen \(AC\), so ist die Aufgabe eine sehr leichte und zusammenfallend
mit der Aufgabe der Pendelschw[ingungen] auf nicht rotirender Erde.
Sobald man aber \(B\) nach gegebenem Gesetz als beweglich betrachtet, so
verhält es sich ganz anders, die Aufgabe ist eine gleichsam specifisch verschiedene
und sehr schwer. Erleichtert wird sie allerdings sehr, wenn man
die Art der Bewegung von \(B\) so annimmt, dass sie gleichförmig in einem
Kreise geschieht, dessen Planum durch eine von \(C\) dagegen gezogene Normale
im Centrum des Kreises getroffen wird. Das erstere ist der Fall, wo der
\HistoricalFoot{XI1.}{6}
\clearpage\SourceLeaf{47}\HistoricalHead{42}{L}{PHYSIK. BRIEFWECHSEL.}
\noindent \textsc{Foucauld}sche Versuch unter dem Pole gemacht wird, das andere unter jeder
anderen Breite. Wer aus dem Umstande, dass im ersteren Falle die Auflösung
so leicht ist, sich zu der Erwartung verleiten lässt, auch im zweiten
(specifisch ganz verschiedenen) Fall eben so leicht oder fast eben so leicht fertig
werden zu können, täuscht sich. Es gibt keinen andern. Rath, als in einer
oder der anderen Form die strenge Behandlung durchzumachen. \textsc{Planas} Abhandlung
habe ich leider allen Suchens ungeachtet noch immer nicht wieder
finden können. Eine andere, im ganzen ähnliche Behandlung der Aufgabe,
gestützt auf meine Grundformeln in \textsc{Benzenbergs} Buch[\textsuperscript{*})]\begingroup\renewcommand{\thefootnote}{}\footnotetext{[\textsuperscript{*}) Siehe \textit{Fundamentalgleichungen für die Bewegung schwerer Körper auf der rotirenden Erde}, Werke V, S. 496.]}\endgroup, hat \textsc{Clausen} in
den Bulletins der Petersburger Akademie gegeben[\textsuperscript{**})]\begingroup\renewcommand{\thefootnote}{}\footnotetext{[\textsuperscript{**}) H. \textsc{Clausen}, \textit{Über den Einfluss der Umdrehung der Gestalt der Erde auf die scheinbare Bewegung an der Oberfläche derselben}, Bulletin phys.-math. de l'Académie de St.-Pétersbourg \textbf{10}, 1852.]}\endgroup, die ich aber auch in
diesem Augenblicke nicht genauer nachweisen kann.\quad .\ .\ .\ .\ .\ .\ .\ .\ .\ .\ .\ .\ .\ .
\clearpage\endgroup
